\section*{Apresentação da entrega}
Esta é a primeira de uma série de entregas incrementais do Projeto Integrador. Utilize as discussões, análises e produções realizadas nas aulas e no mini-workshop como ponto de partida. Todos os itens indicados neste template deverão ser apresentados nesta entrega, ainda que tenham sido inicialmente discutidos ou esboçados em atividades anteriores.

Revise, aprofunde, integre e documente esses resultados de acordo com as orientações de cada seção. Os componentes produzidos nesta entrega poderão ser revisitados e refinados nas próximas etapas do Projeto Integrador, à medida que novos conteúdos forem trabalhados na disciplina.

\subsection*{Relação com a avaliação final do Projeto Integrador}
O projeto final possui 10 quesitos de avaliação, totalizando 10,0 pontos. Esta entrega trabalha apenas os componentes compatíveis com os conteúdos já estudados. A tabela abaixo mostra o mapa completo e evita antecipar etapas futuras.

\begin{table}[!ht]\centering\scriptsize\renewcommand{\arraystretch}{1.35}
\begin{tabularx}{\linewidth}{|>{\centering\arraybackslash}p{0.65cm}|>{\raggedright\arraybackslash}X|>{\centering\arraybackslash}p{1.2cm}|>{\centering\arraybackslash}p{2cm}|>{\raggedright\arraybackslash}X|}
\hline
\textbf{Nº} & \textbf{Quesito da avaliação final} & \textbf{Pts.} & \textbf{Nesta entrega} & \textbf{Orientação} \\
\hline
1 & Problema, contexto e objetivos & 1,0 & Aprofundar & Clareza, delimitação, justificativa e objetivos \\
\hline
2 & Visão sistêmica e sociotécnica & 1,0 & Consolidar & Ambiente, fronteiras, entradas/saídas e articulação sociotécnica \\
\hline
3 & Partes interessadas e necessidades & 1,0 & Depois & Será aprofundado em 29/09--01/10 \\
\hline
4 & Processos organizacionais & 1,0 & Consolidar & Recuperar/refinar o AS-IS já trabalhado \\
\hline
5 & Dados e informação & 1,0 & Consolidar & Registrar o que já foi identificado nas aulas/workshop \\
\hline
6 & Requisitos do SI & 1,5 & Depois & Será trabalhado em 06/10--08/10 \\
\hline
7 & Modelagem conceitual & 1,0 & Depois & Entidades, atributos, relacionamentos e DER \\
\hline
8 & Prototipação e integração & 1,0 & Depois & Interfaces, fluxos e coerência \\
\hline
9 & Privacidade, segurança e governança & 0,75 & Depois & Minimização, dados pessoais, CID e governança \\
\hline
10 & Integração, apresentação e qualidade final & 0,75 & Final & Coerência global, documentação e apresentação \\
\hline
\end{tabularx}
\end{table}
